\section{ReAct 范式回顾}

\subsection{从 CoT 到 ReAct}

ReAct 是 Agent 领域最经典的工作之一（姚顺雨团队，2022 年底 / 2023 年初发表于 NeurIPS，引用量已超 4000）。它的演进脉络如下：

\begin{enumerate}
    \item \textbf{CoT（Chain of Thought）}：面向回答问题——``Let's think step by step''，启发了后续的 Prompt Engineering。
    \item \textbf{Act Only}：语言模型可以和环境交互，产生动作（Action），环境返回观察（Observation），但没有显式思考过程。
    \item \textbf{ReAct = Reasoning + Acting}：在产生动作之前先进行 Reasoning（Thought），降低模型幻觉，更好地 grounding 到合理的动作。
\end{enumerate}

\begin{importantbox}{ReAct 的核心循环}
Thought $\rightarrow$ Action $\rightarrow$ Observation $\rightarrow$ Thought $\rightarrow$ ... $\rightarrow$ Final Answer

ReAct 启发了后续 Agent 的设计和开发，也可能启发了 GPT API 中 Function Calling 特性的诞生（2023年6月）。
\end{importantbox}

\begin{knowledgebox}{关于 Reasoning Model}
现代的 Reasoning Model 在回答问题之前已经学会了思考，不需要再显式地通过 Prompt 告诉它"你要先思考"。后续介绍 Reasoning Model 时会展开讨论其使用规范和标准。
\end{knowledgebox}

\subsection{本章小结}
ReAct 将推理（Reasoning）与行动（Acting）结合，是 Agent 设计的基石。理解其 Thought-Action-Observation 三部曲是理解后续所有 Agent 框架的前提。

